\chapter{PRELIMINARIES}

Introductory paragraph for Chapter 2.

\section{Section 1}

Introductory paragraph for Section 1. 

\begin{defn}\label{ss}
	\cite{BkEx} Let $\R$~\label{real} denote the set of real numbers. A \textit{successor set} is a subset $S$ of $\R$ with the following properties:
	\begin{enumerate}
		\item[(i)] $1\in S$; and
		\item[(ii)] If $n\in S$, then $n+1\in S$.
	\end{enumerate}
\end{defn}

\begin{ex}
	The set $\R^+$~\label{preal} of positive real numbers is a successor set, since it satisfies properties (i) and (ii) of Definition~\ref{ss}. The set $\R^-$~\label{nreal} of negative real numbers is not a successor set, since $1\notin\R^-$. The open interval $(0,10)$ is not a successor set, because $9\in(0,10)$ but $9+1=10\notin(0,10)$.
\end{ex}

\begin{prop}\label{ints}
	The intersection of any collection of successor sets is a successor set.
\end{prop}

\begin{proof}
	Let $\mathcal{S}$ be a collection of successor sets, and set $T:=\ds\bigcap_{S\in\mathcal{S}} S$. We will show that $T$ is a successor set by verifying properties (i) and (ii) of Definition~\ref{ss} for the set $T$.
	\begin{enumerate}
		\item[(i)] Because each set $S\in\mathcal{S}$ is a successor set, $1\in\mathcal{S}$ for all $S\in\mathcal{S}$. Thus, $1\in T$.
		\item[(ii)] Let $n\in T$. Then $n\in S$ for all $S\in\mathcal{S}$. Since each set $S\in\mathcal{S}$ is a successor set, it follows that $n+1\in S$ for all $S\in\mathcal{S}$. Hence, $n+1\in T$.
	\end{enumerate}
	Therefore, $T$ is a successor set.
\end{proof}

\begin{defn}\label{nat}
	\cite{BkEx} The set of \textit{natural numbers}, denoted by $\N$, is defined to be the intersection of all successor sets of $\R$. In symbols,
	\[
	\N:=\bigcap\{S\subseteq\R: S\text{ is a successor set}\}.
	\]
\end{defn}

Proposition~\ref{ints} guarantees that $\N$ is a successor set. In fact, by definition, $\N$ is the ``smallest'' among all successor sets of $\R$ in the sense that $\N\subseteq S$ for all successor sets $S$.

From here, we prove the following result that is used to prove statements involving natural numbers. See the book of \cite{BkEx}.

\begin{thm}[Principle of Mathematical Induction]\label{pmi}
	Let $P_n$ be a statement involving $n\in\N$. Suppose that the following conditions hold:
	\begin{description}
		\item[(PM1)] $P_1$ is true.
		\item[(PM2)] For $k\in\N$, if $P_k$ is true, then $P_{k+1}$ is true.
	\end{description}
	Then $P_n$ is true for all $n\in\N$.
\end{thm}

\begin{proof}
	Let $S:=\{m\in\N: P_m\text{ is true}\}.$ The conclusion of the theorem follows if we can prove that $S=\N$. By definition of $S$, we have $S\subseteq\N$, so it is enough to show that $\N\subseteq S$.
	
	For this, we claim that $S$ is a successor set. Since $P_1$ is true by \textbf{(PM1)}, $1\in S$. Suppose that $n\in S$. Then $P_n$ is true. By \textbf{(PM2)}, $P_{n+1}$ is true, so that $n+1\in S$. Thus, $S$ is a successor set.
	
	Since $\N$ is the intersection of all successor sets, it follows that $\N\subseteq S$. Therefore, $S=\N$.
\end{proof}

\begin{ex}
	Prove that for all $n\in\N$,
	\begin{align}\label{pmiex}
		1^3 + 2^3 + \cdots + n^3 = \frac{n^2(n+1)^2}{4}.
	\end{align}
\end{ex}

\begin{proof}
	For $n\in\N$, let $P_n$ be statement (\ref{pmiex}). Then $P_1$ is the statement
	\begin{align}\label{p1}
		1^3 = \dfrac{1^2(1+1)^2}{4},    
	\end{align}
	which is true because both sides of Equation~(\ref{p1}) are equal to 1. To finish the proof using Theorem~\ref{pmi}, it suffices to prove \textbf{(PM2)} for the statement $P_n$. Suppose that for $k\in\N$, $P_k$ is true, that is,
	\begin{align}\label{pk}
		1^3 + 2^3 + \cdots + k^3 = \frac{k^2(k+1)^2}{4}.
	\end{align}
	We claim that $P_{k+1}$ is true, that is,
	\begin{align*}
		1^3 + 2^3 + \cdots + (k+1)^3 = \frac{(k+1)^2(k+2)^2}{4}.
	\end{align*}
	Indeed, by using (\ref{pk}) and algebraic manipulation, we have
	\begin{align*}
		1^3 + 2^3 + \cdots + (k+1)^3 &= \left(1^3 + 2^3 + \cdots + k^3\right) + (k+1)^3\\
		&= \frac{k^2(k+1)^2}{4} + (k+1)^3\\
		&= \frac{k^2(k+1)^2+4(k+1)^3}{4}\\
		&= \frac{(k+1)^2\big(k^2+4(k+1)\big)}{4}\\
		&= \frac{(k+1)^2\big(k^2+4k+4\big)}{4}\\
		&= \frac{(k+1)^2(k+2)^2}{4}.
	\end{align*}
\end{proof}

\section{Section 2}

Introductory paragraph for Section 2.

\begin{defn}\label{group}
	A \textit{group} is a set $G$ together with a map $*$ from $G\times G$ into $G$ (where for $a,b\in G$, the image of $(a,b)$ under $*$ is denoted by $a*b$\label{binop}), called a \textit{binary operation} on $G$, such that the following hold.
	\begin{description}
		\item[(G1)] The binary operation $*$ is \textit{associative}:
		\[
		(a*b)*c = a*(b*c),\quad\text{for all } a,b,c\in G.
		\]
		\item[(G2)] There exists an element in $G$ called the \textit{identity} in $G$ under $*$, denoted by $e$ (or $e_G$~\label{ident} if $G$ needs to be indicated), such that
		$a*e=e*a=a$ for all $a\in G$.
		\item[(G3)] For each $a\in G$, there exists an element in $G$ called the \textit{inverse} of $a$ in $G$ under $*$, denoted by $a^{-1}$~\label{inv}, such that $a*a^{-1}=a^{-1}*a=e.$
	\end{description}
\end{defn}

\begin{ex}
	Let $G:=\{u,v,w,x,y,z\}$ and define a binary operation $*$ on $G$ as follows: for $a,b\in G$, $a*b$ is defined to be the element in $G$ on Table~\ref{deesix} located in the same row as that of the element $a$ before the vertical line and in the same column as that of the element $b$ above the horizontal line. For instance, according to the table, we have $w*y=x$ and $y*w=v$. The fact that there is exactly one element in $G$ in each entry on the lower right block of the table ensures that $*$ is a well-defined binary operation on $G$. This table is called a \textit{Cayley table} for the binary operation $*$.
	
	\begin{table}[!htbp]
		\centering
		\caption{A Cayley table for a binary operation $*$ on the set $\{u,v,w,x,y,z\}$}
		\vspace{10pt}
		\begin{tabular}{c|c c c c c c} 
			$*$ & $u$ & $v$ & $w$ & $x$ & $y$ & $z$ \\ 
			\hline
			$u$ & $u$ & $v$ & $w$ & $x$ & $y$ & $z$ \\
			$v$ & $v$ & $u$ & $y$ & $z$ & $w$ & $x$ \\
			$w$ & $w$ & $z$ & $u$ & $y$ & $x$ & $v$ \\
			$x$ & $x$ & $y$ & $z$ & $u$ & $v$ & $w$ \\
			$y$ & $y$ & $x$ & $v$ & $w$ & $z$ & $u$ \\
			$z$ & $z$ & $w$ & $x$ & $v$ & $u$ & $y$
		\end{tabular}
		\label{deesix}
	\end{table}
	
	\noindent We will show that $G$ is a group under $*$. Note that $*$ is associative. Indeed, we can verify using Table~\ref{deesix} that
	\[
	(w*x)*z=y*z=u=w*w=w*(x*z),
	\]
	and similar computations hold for all other choices of elements in $G$. Next, the identity element in $G$ under $*$ is $u$, because $v*u=u*v=v$, and similar statements hold for the other elements in $G$. Finally, since $y*z=z*y=u$, we have $y^{-1}=z$ and $z^{-1}=y$, and a quick inspection of Table~\ref{deesix} shows that $u^{-1}=u$, $v^{-1}=v$, $w^{-1}=w$, and $x^{-1}=x$. Therefore, $G$ is a group under $*$.
\end{ex}

\section{Section 3}

Introductory paragraph for Section 3. All definitions are taken from \cite{cl05}.

\begin{defn}\label{gr}
	A \textit{graph} $G$ is a pair $(V(G),E(G))$ consisting of a nonempty finite set $V(G)$ and a set $E(G)$ of unordered pairs of elements of $V(G)$. The sets $V(G)$~\label{vg} and $E(G)$~\label{eg} are called the \textit{vertex set} and the \textit{edge set} of the graph $G$, respectively. The elements of $V(G)$ are called \textit{vertices} in $G$, while the elements of $E(G)$ are called \textit{edges} in $G$. The cardinalities of $V(G)$ and $E(G)$ are called the \textit{order} and the \textit{size} of $G$, respectively.
\end{defn}

Usually, a graph is represented visually by drawing a point or circle for every vertex of the graph, and drawing a curve (usually a line segment) connecting two vertices, say, $v$ and $w$, if the unordered pair $\{v,w\}$ (written briefly as $vw$) is an edge of the graph.

\begin{ex}
	Figure~\ref{graphex} is an illustration of a graph $G$ with vertex set $V(G):=\{v_1, v_2, v_3, v_4, v_5, v_6, v_7\}$
	and edge set $E(G):=\{v_1v_3, v_1v_4, v_2v_4, v_3v_4,\lb v_4v_5, v_4v_6, v_4v_7, v_5v_6\}.$ In this case, the order of $G$ is 7, while its size is 8.
\end{ex}

\begin{figure}[H]
	\[
	\pic
	\Vertex (0,35) (70,0) (70,70) (140,35) (210,0) (210,70) (280,35)
	\Edge (0,35) (140,35)
	\Edge (140,35) (280,35)
	\Edge (0,35) (70,0)
	\Edge (70,70) (140,35)
	\Edge (140,35) (210,0)
	\Edge (210,70) (210,0)
	\Edge (140,35) (210,70)
	\Edge (140,35) (70,0)
	\Align[c] ($v_1$) (-7,35)
	\Align[c] ($v_2$) (62,70)
	\Align[c] ($v_3$) (62,-5)
	\Align[c] ($v_4$) (140,43)
	\Align[c] ($v_5$) (218,70)
	\Align[c] ($v_6$) (218,-5)
	\Align[c] ($v_7$) (288,35)
	\cip
	\]
	\caption{A graph with 7 vertices and 8 edges}
	\label{graphex}
\end{figure}